\chapter{Архитектурное решение после закрытия вихрь--Хиггс коннектора}
\label{ch:v6-spectral-transition-connector-architecture-branch-decision-gate}

\section{Поставленная развилка}

Предыдущие три гейта проверили все минимальные способы связать уже
построенный семейный вихрь $(Q,T,B)$ со сменой ранга хиггс-разрешённой
спектральной опоры. Радиальный квадрат $M_{300}$ не создаёт $H=0$ в
сердцевине; прямые эквивариантные стрелки $B,Q,T\to H$ отсутствуют;
корректное градуированное произведение и центрированная моритова кривизна
не порождают смешанного портала.

После этих отрицательных результатов остаются три логические возможности:
\begin{enumerate}
 \item продолжать искать скрытый коннектор внутри того же родителя;
 \item добавить новый бифундаментальный носитель или высший функционал;
 \item признать спектральный и вихревой объекты двумя самостоятельными
 секторами и продолжить каждый только его собственными средствами.
\end{enumerate}

Настоящая глава является архитектурным решением, а не новой физической
теорией. Она запрещает превращать отрицательный результат в неявное
расширение модели.

\section{Строгий ledger закрытых механизмов}

Пусть $s=|H|^2$ и $T=\operatorname{Tr}(Q^2)$. Для существующего
радиального квадрата точечный минимум равен
\begin{equation}
 s_*(T)=\frac{T}{6}+\frac12\geq\frac12.
\end{equation}
Даже оптимистический единичный ядровый штраф оставляет
$s_*\geq1/4$. Поэтому вихрь не локализует регулярный переход нейтринной
опоры $W_\nu(H)$ из ранга ноль в ранг один.

Представления семейных полей имеют спины $1,2,3$ и являются слабыми
синглетами, тогда как $H$ является слабым дублетом. Отсюда
\begin{equation}
 \operatorname{Hom}_G(B,H)=
 \operatorname{Hom}_G(Q,H)=
 \operatorname{Hom}_G(T,H)=0.
\end{equation}
Существование окружающей полной матричной алгебры не меняет диаграмму
физических стрелок; это соответствует стандартному чтению конечных
спектральных троек \cite{krajewski1998}.

Наконец, для корректного градуированного произведения
\begin{equation}
 D_{\rm tot}=D_F\otimes1+\gamma_F\otimes D_H
\end{equation}
смешанная часть квадрата сокращается из-за
$\gamma_FD_F+D_F\gamma_F=0$. Обычный фактор двухформ удаляет
бесследовую семейную форму, а центрированная относительная кривизна
удовлетворяет
\begin{equation}
 \tau(R_E^2)=\tau(F_F^2)+\tau(F_H^2).
\end{equation}
Это согласуется с градуированным произведением и суперсвязностным
исчислением \cite{quillen1985,kucerovsky1997}.

Три механизма логически независимы: первый проверяет знак потенциала,
второй --- представления степени один, третий --- составную степень два.
Их совместный провал закрывает поиск \emph{скрытого} коннектора в
текущем родителе.

\section{Почему новый портал не является продолжением той же модели}

Инвариант
\begin{equation}
 \lambda_{QH}\operatorname{Tr}(Q^2)H^\dagger H
\end{equation}
разрешён симметриями как эффективный член. Однако существующие след,
кривизна и фактор форм не определяют ни его наличие, ни знак, ни
коэффициент $\lambda_{QH}$. Его добавление поэтому является новой
динамической гипотезой.

Аналогично, смешанный модуль одноформ существует как пространство
возможных коннекторов, но не содержит канонически выбранной ненулевой
секции. Выбор элемента такого модуля меняет оператор Дирака и конечную
диаграмму. По смыслу сильной эквивалентности Мориты сам бимодуль задаёт
тип перехода, но не выделяет произвольную связность без дополнительного
принципа \cite{brown-green-rieffel1977}.

Следовательно, оба ремонта допустимы только после явного объявления новой
версии, повторного аудита следа, аномалий, вакуума, нормировки и полного
гессиана. Внутри Тома VI они не используются.

\section{Выбранная архитектура}

\begin{theorem}[Разделение двух строгих секторов]
В текущей версии спектральный конечранговый дефект и бозонная вихревая
нить сохраняются как два самостоятельных объекта. Их топологические
классы и устойчивость не отменяются отсутствием портала, но проект не
отождествляет вихрь с наблюдаемой частицей и не использует его для
локализации хиггс-разрешённой смены ранга.
\end{theorem}

Спектральная ветвь наследует
\begin{equation}
 N,\quad U,\quad [N,U]=U,\quad
 Q_{\rm def}=|e_0\rangle\langle e_0|\otimes q_0,
\end{equation}
компонентные классы $12+3$, хиггс-разрешённое разложение
$6+6+2+1$ и регулярный квадратичный носитель $W_\nu(H)$.

Бозонная ветвь наследует конечный стационарный профиль $(Q,T,B)$,
положительный поперечный спектр прямой нити и действие Намбу--Гото для
переносных мод. Она может исследоваться как самостоятельный
топологический или скрытый сектор, но такое физическое чтение пока не
является идентификацией с наблюдаемой тёмной материей.

Разделение не означает доказанного полного отсутствия взаимодействий.
Оно означает лишь отсутствие выведенного внутреннего коннектора в
замороженной архитектуре. В частности, универсальная гравитационная
связь не объявляется, пока проект не построил ковариантное действие с
динамической метрикой.

\section{Следующий различающий тест}

После отделения вихря первый допустимый вопрос относится к самой
спектральной ветви. Носитель $W_\nu(H)$ меняет ранг только при $H=0$.
Следующий гейт
\texttt{version6\_spectral\_transition\_higgs\_vacuum\_topology\_localization\_gate}
должен проверить, может ли конфигурационное пространство одного
хиггсовского дублета поддерживать топологически устойчивый локальный нуль
без семейного вихря и без нового поля.

Нулевая гипотеза: вакуумное многообразие нормированного комплексного
дублета не имеет нужной низшей гомотопии, поэтому из одного $H$ нельзя
получить устойчивый точечный, линейный или доменный дефект, локализующий
смену ранга. Классификация дефектов по гомотопии вакуумного многообразия
является стандартной \cite{mermin1979}.

Стоп-критерий: если после точного учёта калибровочного фактора все
соответствующие классы тривиальны, внутренняя топологическая локализация
$W_\nu$ закрывается. Тогда остаются либо нетопологическая динамическая
локализация с выведенным потенциалом, либо явно новая архитектура.

\begin{nogocriterion}[Запрет скрытого переоткрытия портала]
Нельзя в последующих спектральных гейтах использовать профиль $Q,T,B$,
коэффициент $\lambda_{QH}$ или выбранную смешанную одноформу, не объявив
новую модель и не повторив её родительские проверки.
\end{nogocriterion}

\begin{protocol}
\begin{enumerate}
 \item радиальная локализация через $M_{300}$: \textbf{закрыта};
 \item прямой эквивариантный коннектор: \textbf{отсутствует};
 \item двухшаговый моритов коннектор: \textbf{отсутствует};
 \item поиск ещё одного скрытого портала в текущем родителе:
 \textbf{остановлен};
 \item новый бифундаментальный носитель или высший функционал:
 \textbf{только новая явно объявленная модель};
 \item спектральный и вихревой секторы в версии VI:
 \textbf{разделены};
 \item вихрь как наблюдаемая материя: \textbf{не выведен};
 \item следующий тест: \textbf{топология хиггсовского вакуума и
 внутренняя локализация смены ранга};
 \item физическое замыкание: \textbf{не достигнуто}.
\end{enumerate}
\end{protocol}

Машинный ledger решения сохранён в
\path{s2t_v6_spectral_transition_connector_architecture_branch_decision_gate_results.json}.

\begin{thebibliography}{9}
\bibitem{krajewski1998}
T.~Krajewski,
\emph{Classification of Finite Spectral Triples},
Journal of Geometry and Physics 28 (1998), 1--30.

\bibitem{quillen1985}
D.~Quillen,
\emph{Superconnections and the Chern Character},
Topology 24 (1985), 89--95.

\bibitem{kucerovsky1997}
D.~Kucerovsky,
\emph{The KK-Product of Unbounded Modules},
K-Theory 11 (1997), 17--34.

\bibitem{brown-green-rieffel1977}
L.~G.~Brown, P.~Green, M.~A.~Rieffel,
\emph{Stable Isomorphism and Strong Morita Equivalence of
$C^*$-Algebras},
Pacific Journal of Mathematics 71 (1977), 349--363.

\bibitem{mermin1979}
N.~D.~Mermin,
\emph{The Topological Theory of Defects in Ordered Media},
Reviews of Modern Physics 51 (1979), 591--648.
\end{thebibliography}